\documentclass[12pt,a4paper]{article}

% ============================================================
% PACKAGES
% ============================================================

\usepackage[utf8]{inputenc}
\usepackage[T1]{fontenc}
\usepackage[french]{babel}

\usepackage{geometry}
\geometry{
    left=2.5cm,
    right=2.5cm,
    top=2.5cm,
    bottom=2.5cm
}

\usepackage{graphicx}
\usepackage{float}
\usepackage{amsmath}
\usepackage{amssymb}
\usepackage{booktabs}
\usepackage{array}
\usepackage{xcolor}
\usepackage{listings}
\usepackage{caption}
\usepackage{subcaption}
\usepackage{hyperref}
\usepackage{fancyhdr}
\usepackage{enumitem}
\usepackage{setspace}

% ============================================================
% CONFIGURATION PYTHON
% ============================================================

\definecolor{codegreen}{rgb}{0,0.5,0}
\definecolor{codegray}{rgb}{0.45,0.45,0.45}
\definecolor{codepurple}{rgb}{0.58,0,0.82}
\definecolor{backcolour}{rgb}{0.96,0.96,0.96}

\lstdefinestyle{python}{
    backgroundcolor=\color{backcolour},
    commentstyle=\color{codegreen},
    keywordstyle=\color{blue},
    numberstyle=\tiny\color{codegray},
    stringstyle=\color{codepurple},
    basicstyle=\ttfamily\small,
    breakatwhitespace=false,
    breaklines=true,
    captionpos=b,
    keepspaces=true,
    numbers=left,
    numbersep=5pt,
    showspaces=false,
    showstringspaces=false,
    showtabs=false,
    tabsize=4,
    language=Python
}

\lstset{style=python}

% ============================================================
% EN-TÊTES
% ============================================================

\pagestyle{fancy}
\fancyhf{}
\fancyhead[L]{TP1 -- Statistiques descriptives univariées}
\fancyhead[R]{CIPRIANI -- VILLENEUVE}
\fancyfoot[C]{\thepage}

\setlength{\headheight}{15pt}

% ============================================================
% DOCUMENT
% ============================================================

\begin{document}

% ============================================================
% PAGE DE GARDE
% ============================================================

\begin{titlepage}

    \centering

    \vspace*{2cm}

    \vspace{2cm}

    {\Huge\bfseries
    TP1\\[0.3cm]
    Statistiques descriptives univariées
    \par}

    \vspace{1cm}

    {\Large
    Analyse du jeu de données Iris
    \par}

    \vfill

    \begin{tabular}{rl}
        \textbf{Étudiants :} & CIPRIANI -- VILLENEUVE \\
        \textbf{Date :} & 7 septembre 2026
    \end{tabular}

    \vfill

\end{titlepage}

\tableofcontents
\newpage

% ============================================================
% INTRODUCTION
% ============================================================

% ============================================================
% 2.1
% ============================================================

\section{Analyse préalable}

\subsection{Chargement des données -- Question 2.1}

Le fichier contenant les données est chargé à l'aide de la bibliothèque
\texttt{NumPy}.

Le code utilisé est le suivant :

\begin{lstlisting}
from pprint import pprint
import numpy as np

data = np.genfromtxt(
    "iris.csv",
    delimiter=",",
    names=True,
    dtype=None,
    encoding="utf-8"
)

sepallength = data["sepallength"]
sepalwidth = data["sepalwidth"]
petallength = data["petallength"]
petalwidth = data["petalwidth"]
species = data["class"]

speciesname = np.unique(species)
\end{lstlisting}

La fonction \texttt{np.genfromtxt()} permet de lire le fichier CSV et d'en extraire les données grâce à un délimiteur de données (ici le ",").

L'instruction :

\begin{lstlisting}
speciesname = np.unique(species)
\end{lstlisting}

permet quant à elle de récupérer toutes les valeurs "unique" prises
par la variable \texttt{species}:

\begin{lstlisting}
Species : ['Iris-setosa' 'Iris-versicolor' 'Iris-virginica']
\end{lstlisting}

Il est possible de faire un parallèle avec la fonction python \texttt{Counter()} qui compte et a pour sortie  un dictionnaire avec le nombre de mots uniques:
\begin{lstlisting}
liste = [
    "Iris-setosa",
    "Iris-setosa",
    "Iris-setosa",
    "Iris-versicolor",
    "Iris-versicolor",
    "Iris-virginica"
]
compteur = Counter(liste)
\end{lstlisting}
La sortie est la suivante :
\begin{lstlisting}
Counter({
    "Iris-setosa": 3,
    "Iris-versicolor": 2,
    "Iris-virginica": 1
})
\end{lstlisting}


Dans le jeu de données étudié, celles-ci sont :

\begin{itemize}
    \item Iris-setosa ;
    \item Iris-versicolor ;
    \item Iris-virginica.
\end{itemize}

% ============================================================
% QUESTION 2.2
% ============================================================

\subsection{Nombre d'individus statistiques -- Question 2.2}

Le nombre d'individus est déterminé à partir du nombre de lignes
présentes dans le tableau :

\begin{lstlisting}
nb_individu = len(data)

print(
    f"nombre d'individus statistiques : {nb_individu}"
)
\end{lstlisting}

On obtient :

\[
\boxed{n = 150}
\]

Le jeu de données contient donc \textbf{150 individus statistiques}.


% ============================================================
% QUESTION 2.3
% ============================================================

\subsection{Variables qualitatives -- Question 2.3}

Le jeu de données contient une variable qualitative :

\[
\boxed{\texttt{class}}
\]

également appelée \texttt{species} dans notre programme.

Cette variable peut prendre trois types comme présenté précédemment :

\begin{lstlisting}
Species : ['Iris-setosa' 'Iris-versicolor' 'Iris-virginica']
\end{lstlisting}

Cette variable est \textbf{qualitative nominale}, ici l'ordre ne compte pas (donc elle ne peut pas être ordinale).

% ============================================================
% QUESTION 2.4
% ============================================================

\subsection{Variables quantitatives -- Question 2.4}

Les variables quantitatives sont :

\begin{center}
\begin{tabular}{ll}
\toprule
\textbf{Variable} & \textbf{Signification} \\
\midrule
\texttt{sepallength} & Longueur du sépale \\
\texttt{sepalwidth} & Largeur du sépale \\
\texttt{petallength} & Longueur du pétale \\
\texttt{petalwidth} & Largeur du pétale \\
\bottomrule
\end{tabular}
\end{center}

Ces quatre variables représentent des mesures de "distance" (longueur et largeur plus précisément).

Elles sont donc considérées comme des
\textbf{variables quantitatives continues}.


% ============================================================
% 2.2 SPECIES
% ============================================================

\section{Étude de la variable \texttt{species}}

\subsection{Effectifs -- Question 2.5}

Pour déterminer l'effectif associé à chacune des espèces, le programme
suivant est utilisé :

\begin{lstlisting}
dico_data_effectives = dict()

for type in speciesname:
    dico_data_effectives[type] = np.sum(
        species == type
    )

pprint(dico_data_effectives)
\end{lstlisting}
C'est le principe du dictionnaire qui a été utilisé, car plus simple à manipuler puisqu'il associe un label avec des données. En sortie, il est possible d'observer les résultats suivants:

\begin{lstlisting}
{   np.str_('Iris-setosa'): np.int64(50),
    np.str_('Iris-versicolor'): np.int64(50),
    np.str_('Iris-virginica'): np.int64(50)}
\end{lstlisting}


Il est donc possible d'établir le tableau suivant:

\begin{center}
\begin{tabular}{lcc}
\toprule
\textbf{Espèce} & \textbf{Effectif} & \textbf{Fréquence} \\
\midrule
Iris-setosa     & 50 & 33,3 \% \\
Iris-versicolor & 50 & 33,3 \% \\
Iris-virginica  & 50 & 33,3 \% \\
\midrule
Total           & 150 & 100 \% \\
\bottomrule
\end{tabular}
\end{center}

Les trois espèces possèdent donc exactement le même effectif.

Le jeu de données est ainsi parfaitement équilibré par rapport à la
variable \texttt{species}.

% ============================================================
% QUESTION 2.6
% ============================================================

\subsection{Représentations graphiques -- Question 2.6}

Deux représentations graphiques ont été utilisées pour visualiser la
répartition des espèces :

\begin{itemize}
    \item un diagramme en secteurs ;
    \item un histogramme.
\end{itemize}

Le code utilisé est :

\begin{lstlisting}
import matplotlib.pyplot as plt

labels = dico_data_effectives.keys()
sizes = dico_data_effectives.values()

fig, (ax1, ax2) = plt.subplots(
    1, 2, figsize=(10, 4)
)

ax1.pie(
    sizes,
    labels=labels
)

ax2.hist(
    species,
    bins=3,
    edgecolor="black",
    label=labels
)

plt.show()
\end{lstlisting}

\begin{figure}[H]
    \centering
    \includegraphics[width=1\textwidth]{captures/graphique_2_2.png}
    \caption{Répartition des individus selon l'espèce.}
    \label{fig:species}
\end{figure}

Le diagramme en secteurs montre une répartition parfaite entre les différentes \texttt{species}.

L'histogramme confirme les résultats précédents puisque les trois barres possèdent une hauteur identique de 50 individus.


% ============================================================
% 2.3 PETALLENGTH
% ============================================================

\section{Étude de la variable \texttt{petalLength}}


% ============================================================
% QUESTION 2.7
% ============================================================

\subsection{Histogrammes -- Question 2.7}

L'histogramme permet de représenter la distribution d'une variable
quantitative continue en regroupant les observations dans différents
intervalles.

Le code général utilisé est :

\begin{lstlisting}
ax1.hist(
    petallength,
    bins=intervalle,
    edgecolor="black"
)
\end{lstlisting}

Un histogramme cumulatif est également réalisé avec :

\begin{lstlisting}
ax2.hist(
    petallength,
    bins=intervalle,
    cumulative=True,
    edgecolor="black"
)
\end{lstlisting}

Pour étudier l'influence du nombre d'intervalles sur la
représentation, quatre valeurs ont été testées :

\[
5,\qquad10,\qquad15,\qquad20.
\]

% ============================================================
% HISTOGRAMMES
% ============================================================

\subsubsection{Histogramme avec 5 intervalles}

\begin{figure}[H]
    \centering
    \includegraphics[width=0.78\textwidth]
    {captures/histogramme_petallength_5_intervalles.png}
    \caption{Histogramme de \texttt{petalLength} avec 5 intervalles.}
    \label{fig:hist5}
\end{figure}

Avec seulement cinq intervalles, la représentation est relativement
grossière.

On distingue néanmoins deux grandes zones :

\begin{itemize}
    \item un premier groupe autour de 1 à 2 cm ;
    \item un second groupe beaucoup plus étendu, approximativement
    entre 3 et 7 cm.
\end{itemize}

Une zone de très faible effectif apparaît entre ces deux groupes.

% ============================================================

\subsubsection{Histogramme avec 10 intervalles}

\begin{figure}[H]
    \centering
    \includegraphics[width=0.78\textwidth]
    {captures/histogramme_petallength_10_intervalles.png}
    \caption{Histogramme de \texttt{petalLength} avec 10 intervalles.}
    \label{fig:hist10}
\end{figure}

Avec dix intervalles, la structure de la distribution devient plus
précise.

Le premier groupe est fortement concentré autour des petites longueurs
de pétales.

Le deuxième groupe est plus étalé et possède son maximum d'effectif
aux alentours de 4 à 5 cm.

La séparation entre le premier groupe et le reste des observations
apparaît également plus clairement.

% ============================================================

\subsubsection{Histogramme avec 15 intervalles}

\begin{figure}[H]
    \centering
    \includegraphics[width=0.78\textwidth]
    {captures/histogramme_petallength_15_intervalles.png}
    \caption{Histogramme de \texttt{petalLength} avec 15 intervalles.}
    \label{fig:hist15}
\end{figure}

Avec quinze intervalles, davantage de détails apparaissent dans la
distribution.

La concentration du premier groupe autour de 1,5 cm devient
particulièrement visible.

Pour les longueurs supérieures à environ 3 cm, la distribution est
plus étendue et les variations locales des effectifs deviennent plus
faciles à observer.

% ============================================================

\subsubsection{Histogramme avec 20 intervalles}

\begin{figure}[H]
    \centering
    \includegraphics[width=0.78\textwidth]
    {captures/histogramme_petallength_20_intervalles.png}
    \caption{Histogramme de \texttt{petalLength} avec 20 intervalles.}
    \label{fig:hist20}
\end{figure}

Avec vingt intervalles, la résolution de l'histogramme est encore
augmentée.

Les fluctuations locales des données deviennent davantage visibles.
Cependant, la représentation devient également moins synthétique :
chaque intervalle contient moins d'observations et la forme générale
de la distribution est donc plus irrégulière.

% ============================================================
% QUESTION 2.8
% ============================================================

\subsection{Analyse de l'influence du nombre d'intervalles -- Question 2.8}

La comparaison des quatre histogrammes met en évidence l'importance du
choix du nombre d'intervalles.

Avec un faible nombre d'intervalles, par exemple cinq, la représentation
est très grossière et ne permet pas d'identifier une partie des
variations internes.

Lorsque le nombre d'intervalles augmente, la largeur de chaque classe
diminue. La représentation devient donc plus détaillée.

Cependant, un nombre trop important d'intervalles peut donner une
impression de forte irrégularité, car chaque classe contient alors un
nombre plus faible d'observations.

Malgré ces différences de représentation, une caractéristique reste
présente quel que soit le nombre d'intervalles : la distribution de
\texttt{petalLength} n'est pas unimodale.

On observe une forte concentration des observations pour des longueurs
comprises approximativement entre 1 et 2 cm, puis une seconde zone
beaucoup plus large entre environ 3 et 7 cm.

La zone comprise approximativement entre 2 et 3 cm contient très peu
d'observations.

Cette structure montre que la population totale n'est pas homogène du
point de vue de la longueur des pétales.

Le choix du nombre d'intervalles ne modifie donc pas les données
elles-mêmes, mais modifie la manière dont leur structure est rendue
visible.

Dans le cas présent, une valeur intermédiaire, par exemple entre
10 et 15 intervalles, constitue un compromis intéressant entre
lisibilité et niveau de détail.

% ============================================================
% STURGES / YULE
% ============================================================

\subsubsection{Choix automatique du nombre d'intervalles}

Le programme contient également deux estimations du nombre
d'intervalles :

\begin{lstlisting}
n = len(petallength)

intervalle_Sturge = int(
    np.ceil(1 + 3.3 * np.log10(n))
)

intervalle_yule = int(
    np.ceil(2.5 * n**0.25)
)
\end{lstlisting}

Pour :

\[
n=150,
\]

la règle de Sturges donne approximativement :

\[
k =
\left\lceil
1+3,3\log_{10}(150)
\right\rceil
\]

soit :

\[
\boxed{k=9}.
\]

La formule utilisée dans le programme pour Yule donne :

\[
k =
\left\lceil
2,5\times150^{1/4}
\right\rceil
\]

soit également environ :

\[
\boxed{k=9}.
\]

Ces valeurs sont proches du choix de dix intervalles testé
précédemment.
Il est aussi à noter qu'il existe un paramètre modifiable \texttt{Bins} qui peut être initialisé à "auto" pour les intervalles. Lors de la simulation, celui-ci renvoie 9 intervalles sur l'histogramme.

% ============================================================
% HISTOGRAMME CUMULATIF
% ============================================================

\subsection{Histogramme cumulatif}

L'histogramme cumulatif représente non plus l'effectif de chaque
intervalle pris indépendamment, mais l'effectif cumulé depuis la plus
petite valeur.

Ainsi, la hauteur de la dernière barre correspond nécessairement au
nombre total d'observations :

\[
150
\]

Cette représentation permet notamment d'observer la proportion
d'individus dont la longueur de pétale est inférieure à une valeur
donnée.

% ============================================================
% QUESTION 2.9
% ============================================================

\subsection{Boîte à moustaches -- Question 2.9}

La boîte à moustaches est construite avec :

\begin{lstlisting}
ax3.boxplot(petallength)
ax3.set_title(
    "Boite a moustache : petallength"
)
\end{lstlisting}

L'histogramme, l'histogramme cumulatif et la boîte à moustaches sont
présentés sur la figure suivante.

\begin{figure}[H]
    \centering
    \includegraphics[width=1\textwidth]
    {captures/graphique_2_3_interv_.png}
    \caption{
        Histogramme, histogramme cumulatif et boîte à moustaches
        de la variable \texttt{petalLength}.
    }
    \label{fig:boxplot}
\end{figure}

Une boîte à moustaches est principalement constituée des éléments
suivants :

\begin{itemize}

    \item le premier quartile \(Q_1\), correspondant au bord inférieur
    de la boîte ;

    \item la médiane \(Q_2\), représentée par le trait situé à
    l'intérieur de la boîte ;

    \item le troisième quartile \(Q_3\), correspondant au bord
    supérieur de la boîte ;

    \item l'écart interquartile :
    \[
    IQR = Q_3-Q_1 ;
    \]

    \item les moustaches, qui s'étendent vers les valeurs extrêmes
    suivant la convention utilisée par le logiciel ;

    \item les éventuelles valeurs aberrantes, qui peuvent être
    représentées individuellement au-delà des moustaches.

\end{itemize}

Dans notre cas :

\[
Q_1 = 1,6\text{ cm}
\]

et :

\[
Q_3 = 5,1\text{ cm}.
\]

L'écart interquartile vaut donc :

\[
IQR = 5,1-1,6
\]

soit :

\[
\boxed{IQR=3,5\text{ cm}}.
\]

La médiane, située à :

\[
\boxed{4,35\text{ cm}},
\]

est également visible à l'intérieur de la boîte.

La boîte est relativement étendue, ce qui traduit une dispersion
importante des longueurs de pétales lorsque les trois espèces sont
considérées simultanément.

% ============================================================
% QUESTION 2.10
% ============================================================

\section{Résumés numériques de \texttt{petalLength}}

\subsection{Calcul des indicateurs -- Question 2.10}

Les différents indicateurs statistiques sont calculés avec NumPy.

Le programme utilisé est :

\begin{lstlisting}
mean_petallength = np.mean(petallength)
median_petallength = np.median(petallength)

ecartType_petallength = np.std(petallength)
variance_petallength = np.var(petallength)

quartile_petallength = np.quantile(
    petallength,
    q=[0.25, 0.75]
)
\end{lstlisting}

Les résultats obtenus par le programme sont :

\begin{center}

\begin{tabular}{lc}
\toprule
\textbf{Indicateur} & \textbf{Valeur} \\
\midrule

Moyenne
& \(3,7587\) cm \\

Médiane
& \(4,35\) cm \\

Écart-type (\texttt{ddof=0})
& \(1,7585\) cm \\

Variance (\texttt{ddof=0})
& \(3,0924\) cm$^2$ \\

Premier quartile \(Q_1\)
& \(1,6\) cm \\

Troisième quartile \(Q_3\)
& \(5,1\) cm \\

Écart interquartile
& \(3,5\) cm \\

\bottomrule
\end{tabular}

\end{center}

\section{ANNEXES}
\subsection{Programme complet}

\begin{lstlisting}
# CIPRIANI VILLENEUVE

# Question 2.1: 
from pprint import pprint
import numpy as np
# Lire le fichier CSV avec numpy (assume un header)
data = np.genfromtxt("iris.csv", delimiter=",",
names=True, dtype=None, encoding="utf-8")
# Colonnes
sepallength = data["sepallength"]
sepalwidth = data["sepalwidth"]
petallength = data["petallength"]
petalwidth = data["petalwidth"]
species = data["class"]
speciesname = np.unique(species)

print(f"Species : {speciesname}")

save_hist_variants: bool = True
status : bool = True #Permet de skip les graphiques

# 2.1 Analyse préalable

# Question 2.2:
nb_individu : int  = len(data)
print(f"nombre d’individus statistiques : {nb_individu}")

# 2.2 Etude de la variable species

# Question 2.5:
dico_data_effectives = dict()
for type in speciesname:
    # dico_data_effectives[type] = {}
    dico_data_effectives[type]= np.sum(species == type)

pprint(dico_data_effectives, indent=4, width=80)

# Question 2.6:
import matplotlib.pyplot as plt

labels = dico_data_effectives.keys()
sizes = dico_data_effectives.values()
print(sizes)

fig, (ax1, ax2) = plt.subplots(1, 2, figsize=(10, 4))
ax1.pie(sizes, labels=labels)
ax1.set_title("Diagramme en secteurs : species")

ax2.hist(species, bins=3, edgecolor='black',label=labels)
ax2.set_title("Histogramme : species")
ax2.set_xlabel("Type de species")
ax2.set_ylabel("Effectif")
# plt.savefig("graphique_2_2.png")
plt.show(block = status)

# 2.3 Étude de la variable petalLength

# Question 2.7:
fig, (ax1, ax2, ax3) = plt.subplots(1, 3, figsize=(17, 7))

# intervalle = 10
intervalle = "auto"
n  = len(petallength)
intervalle_Sturge = int(np.ceil(1 + 3.3 * np.log10(n)))
intervalle_yule = int(np.ceil(2.5 * n**0.25))
print(f"Intervalle Sturge : {intervalle_Sturge}\nIntervalle Yule : {intervalle_yule}")

plt.title(f"Etude de la variable petalLength avec intervalle : {intervalle}")
# ax1.hist(petallength, bins=intervalle_Sturge, edgecolor='black')
ax1.hist(petallength, bins=intervalle, edgecolor='black')
ax1.set_title("Histogramme : petallength")
ax1.set_xlabel("Longueur des petales (cm)")
ax1.set_ylabel("Effectif")
# ax2.hist(petallength, bins=intervalle_Sturge,cumulative = True, edgecolor='black')
ax2.hist(petallength, bins=intervalle,cumulative = True, edgecolor='black')
ax2.set_title("Histogramme cumulatif : petallength")
ax2.set_xlabel("Longueur des petales (cm)")
ax2.set_ylabel("Effectif")

#Dedie au rapport
if save_hist_variants:

    liste_intervalles = [5, 10, 15, 20]

    for nb_intervalles in liste_intervalles:

        fig_hist, ax_hist = plt.subplots(figsize=(8, 6))

        ax_hist.hist(
            petallength,
            bins=nb_intervalles,
            edgecolor="black"
        )

        ax_hist.set_title(
            f"Histogramme petalLength - {nb_intervalles} intervalles"
        )
        ax_hist.set_xlabel("Longueur des pétales (cm)")
        ax_hist.set_ylabel("Effectif")

        nom_fichier = f"histogramme_petallength_{nb_intervalles}_intervalles.png"

        fig_hist.savefig(nom_fichier, dpi=300, bbox_inches="tight")

        plt.close(fig_hist)

# Question 2.8: 

# Question 2.9: 
ax3.boxplot(petallength)
ax3.set_title("Boite à moustache : petallength")
ax3.set_ylabel("Longueur des petales (cm)")
name = "graphique_2_3_interv_"+ str(intervalle) +".png"
# plt.savefig("graphique_2_3_interv_.png")
plt.show(block = status)


# Deuxieme approche : resumes numeriques
# Question 2.10:
mean_petallength = np.mean(petallength)
median_petallength = np.median(petallength)
print(f"Mean : {mean_petallength}\nMedian : {median_petallength}")

ecartType_petallength = np.std(petallength)
variance_petallength = np.var(petallength)
quartile_petallength = np.quantile(petallength, q=[0.25, 0.75])
print(f"Ecart_type : {ecartType_petallength}\nVariance : {variance_petallength}\nQuartile : {quartile_petallength}")

# Mean : 3.758666666666666
# Median : 4.35
# Ecart_type : 1.7585291834055212
# Variance : 3.092424888888889
# Quartile : [1.6 5.1]


\end{lstlisting}


\end{document}